% !TEX program = xelatex
\documentclass[11pt,a4paper]{article}
\usepackage[top=16mm,bottom=15mm,left=19mm,right=19mm]{geometry}
\usepackage{fontspec}
\usepackage{xeCJK}
\usepackage{xcolor}
\usepackage{enumitem}
\usepackage{amsmath}
\usepackage{hyperref}

\setmainfont[
  Path={arial/}, Extension=.ttf,
  UprightFont=ArialMdm, BoldFont=arialceb,
  ItalicFont=ArialMdmItl, BoldItalicFont=ArialCEBoldItalic
]{ArialMdm}
\setCJKmainfont[
  Path={10_SourceHanSansTC/OTF/TraditionalChinese/}, Extension=.otf,
  BoldFont=SourceHanSansTC-Bold
]{SourceHanSansTC-Regular}

\definecolor{navy}{HTML}{183348}
\definecolor{teal}{HTML}{167C80}
\definecolor{ink}{HTML}{263642}
\definecolor{muted}{HTML}{3F4F5B}
\definecolor{rulecolor}{HTML}{CCDCE2}
\hypersetup{colorlinks=true,urlcolor=teal,
  pdftitle={自傳｜楊子賢},pdfauthor={楊子賢},pdfsubject={自傳}}

\pagestyle{empty}
\setlength{\parindent}{0pt}
\setlength{\parskip}{0pt}
\setlength{\emergencystretch}{2em}
\linespread{1.06}
\XeTeXlinebreaklocale "zh"
\XeTeXlinebreakskip = 0pt plus 1pt minus 0.1pt

\newcommand{\autosection}[2]{%
  \vspace{11pt}%
  {\color{navy}\fontsize{13}{16.5}\selectfont\bfseries #1\hfill
  {\color{teal}\fontsize{7.6}{10}\selectfont #2}}\par
  \vspace{3.5pt}{\color{rulecolor}\hrule height 0.7pt}\vspace{8pt}}

\begin{document}
\color{ink}
\fontsize{11}{16.6}\selectfont

%% ---------- 標題 ----------
\begin{center}
  {\color{navy}\fontsize{23}{28}\selectfont\bfseries 自\hspace{0.35em}傳}\par
  \vspace{6pt}
  {\color{muted}\fontsize{9.4}{12}\selectfont
  楊子賢\;ZIH-SIAN YANG｜國立臺灣海洋大學 資訊工程學系｜\href{mailto:jihuty@gmail.com}{jihuty@gmail.com}}
\end{center}
\vspace{2pt}

%% ---------- 求學經歷 ----------
\autosection{求學經歷與學習歷程}{EDUCATION}

我自高中起接觸 C++、Python 程式語言，投入演算法的學習；這段經歷肯定了我的學習成果，也讓我確定了未來的方向。進入大學後我持續參與程式競賽，取得 Kaggle Expert、NCPC 佳作、TOPC 銀獎、CPE 五題與 TOEIC 815 等成績。

在專業課程上，我以演算法與系統兩條主線累積基礎。演算法方面修習\textbf{資料結構}（87／實習 99）、\textbf{演算法}（89）、\textbf{圖論演算法}（91）建立分析問題與證明的習慣；系統方面修習\textbf{作業系統}（86）、\textbf{計算機結構}（87）、\textbf{編譯器}（88）、\textbf{系統程式}（85）、\textbf{嵌入式系統設計}（94）與\textbf{微處理機概論}（91／實驗 97），理解程式在真實硬體上的行為。此外也透過\textbf{人工智慧}、\textbf{機器視覺理論與應用}（88）、\textbf{Python 程式語言}（91）、\textbf{JAVA 程式設計}（98）與\textbf{網頁程式設計}（94）補足實作能力。
%% ---------- 專題研究 ----------
\autosection{專題研究}{RESEARCH}

我與教授共同進行的文獻探討，主題是不可分割物品的公平分配（Fair Division of Indivisible Goods）：當物品相對於總價值變得「微小」時，EF1（Envy-Free up to one item）分配所能保證的納什社會福利（Nash Social Welfare, NSW）能提升到什麼程度。

在相同可加性估值（identical additive valuations）下，已知 EF1 分配的 NSW 至少是最佳解的 $e^{-1/e}\approx 0.692$ 倍。我們探討的定理在物品滿足 $\epsilon$-小物品條件時，給出更強的下界 $\max\{e^{-1/e},\,\rho_n(\epsilon)\}$。其中 $\rho_n(\epsilon)$ 刻劃了 EF1 與小物品限制下最不平均的價值分佈：有 $k$ 人拿到略低於平均的 $\left(1-\frac{k\epsilon}{n}\right)$ 份額，其餘 $n-k$ 人拿到略高於平均的 $\left(1+\frac{(n-k)\epsilon}{n}\right)$ 份額，由此推得保證值隨 $\epsilon$ 變小而趨近最佳分配。

這份研究讓我學會如何閱讀論文、拆解證明結構，並將直覺轉化為可驗證的數學敘述，也讓我確認自己對理論與演算法設計的興趣。

%% ---------- 個人特質與展望 ----------
\autosection{個人特質與未來展望}{OUTLOOK}

長期的競賽訓練養成我面對未知問題時的耐性：先界定問題、再尋找可行解、最後改進效率。我習慣把想法寫成程式驗證，也樂於與同學討論不同的解法。未來我希望在研究所深化機器學習與演算法的訓練，將理論分析與系統實作結合，成為能獨立思考並解決實際問題的研究者。

\end{document}
